% GENERATED FILE. Edit canonical lessons, then run npm run generate:books.
% canonical-source-hash: fnv1a64:fc83a8f65bb1ece2
% canonical-lessons: ES-C654-despegar, ES-C654-aterrizar, ES-C654-piloto, ES-C654-tripulacion, ES-C654-azafata

\chapter{En el avión}
\label{ch:es-en-el-avion}
\hlchaptermodality{654}

\begin{chapteropening}
\textbf{By the end of this chapter:} \emph{I can talk about a flight: take-off and landing, the pilot and the crew.}

Complete the last lesson of chapter 654: I can talk about a flight: take-off and landing, the pilot and the crew.
\end{chapteropening}

\section[despegar]{despegar — to take off; to unstick}

\label{lesson:ES-C654-despegar}



\begin{quote}
Before the new one: say the Spanish for the tray, then the Spanish for the stopover.
\end{quote}

\subsection*{You'll want to know: despegar}
\textbf{despegar} — \textquotedblleft{}to take off; to unstick\textquotedblright{}.

\begin{quote}
El avión despega a las seis.
\end{quote}

\begin{quote}
\emph{The plane takes off at six.}
\end{quote}

The noun is \emph{el despegue}, take-off.

\begin{cousinweb}
\emph{despegar} is built on \emph{pegar} \textquotedblleft{}to stick\textquotedblright{} with \emph{des-} "un-"; \emph{pegar} is inherited from Latin \emph{picāre} \textquotedblleft{}to smear with pitch\textquotedblright{} (EtymDB). A plane unsticks itself from the ground.
\end{cousinweb}

\begin{grammarlens}[title={What you've built}]
One more word for telling, explaining and describing at B1.
\end{grammarlens}

\subsection*{Your turn}
\begin{itemize}
\raggedright
  \item \emph{Say it:} \emph{despegar}
  \item \emph{Say it:} \emph{despegar}, once more
  \item \emph{Say it:} a sentence of your own with \emph{despegar}
  \item \emph{Recall:} say the Spanish for the tray, then the Spanish for the stopover, then say \emph{despegar} again
\end{itemize}

\begin{tcolorbox}[breakable,colback=teal!4,colframe=teal!35!black,title={Before you move on}]
What does despegar mean? (\textquotedblleft{}To take off; to unstick\textquotedblright{}.) Say it once more.
\end{tcolorbox}
\section[aterrizar]{aterrizar — to land}

\label{lesson:ES-C654-aterrizar}



\begin{quote}
Before the new one: say the Spanish for the stopover, then the Spanish for to take off.
\end{quote}

\subsection*{You'll want to know: aterrizar}
\textbf{aterrizar} — \textquotedblleft{}to land\textquotedblright{}.

\begin{quote}
Vamos a aterrizar en diez minutos.
\end{quote}

\begin{quote}
\emph{We are going to land in ten minutes.}
\end{quote}

The noun is \emph{el aterrizaje}. The z becomes c before e: \emph{aterricé}.

\begin{cousinweb}
\emph{aterrizar} is built on \emph{tierra} \textquotedblleft{}land\textquotedblright{}, which is inherited from Latin \emph{terra} (EtymDB), perhaps after French \emph{atterrir} \textquotedblleft{}to land\textquotedblright{}.
\end{cousinweb}

\begin{grammarlens}[title={What you've built}]
One more word for telling, explaining and describing at B1.
\end{grammarlens}

\subsection*{Your turn}
\begin{itemize}
\raggedright
  \item \emph{Say it:} \emph{aterrizar}
  \item \emph{Say it:} \emph{aterrizar}, once more
  \item \emph{Say it:} a sentence of your own with \emph{aterrizar}
  \item \emph{Recall:} say the Spanish for the stopover, then the Spanish for to take off, then say \emph{aterrizar} again
\end{itemize}

\begin{tcolorbox}[breakable,colback=teal!4,colframe=teal!35!black,title={Before you move on}]
What does aterrizar mean? (\textquotedblleft{}To land\textquotedblright{}.) Say it once more.
\end{tcolorbox}
\section[el piloto]{el piloto — the pilot; the driver (racing)}

\label{lesson:ES-C654-piloto}



\begin{quote}
Before the new one: say the Spanish for to take off, then the Spanish for to land.
\end{quote}

\subsection*{You'll want to know: el piloto}
\textbf{el piloto} — \textquotedblleft{}the pilot; the driver (racing)\textquotedblright{}.

\begin{quote}
El piloto habla por el altavoz.
\end{quote}

\begin{quote}
\emph{The pilot is speaking over the loudspeaker.}
\end{quote}

For a woman, \emph{la piloto}.

\begin{cousinweb}
\emph{piloto} was borrowed from Italian \emph{piloto} (EtymDB), first the man who steered a ship into port. Its deeper origin is not settled.
\end{cousinweb}

\begin{grammarlens}[title={What you've built}]
One more word for telling, explaining and describing at B1.
\end{grammarlens}

\subsection*{Your turn}
\begin{itemize}
\raggedright
  \item \emph{Say it:} \emph{el piloto}
  \item \emph{Say it:} \emph{el piloto}, once more
  \item \emph{Say it:} a sentence of your own with \emph{el piloto}
  \item \emph{Recall:} say the Spanish for to take off, then the Spanish for to land, then say \emph{el piloto} again
\end{itemize}

\begin{tcolorbox}[breakable,colback=teal!4,colframe=teal!35!black,title={Before you move on}]
What does el piloto mean? (\textquotedblleft{}The pilot; the driver (racing)\textquotedblright{}.) Say it once more.
\end{tcolorbox}
\section[la tripulación]{la tripulación — the crew}

\label{lesson:ES-C654-tripulacion}



\begin{quote}
Before the new one: say the Spanish for to land, then the Spanish for the pilot.
\end{quote}

\subsection*{You'll want to know: la tripulación}
\textbf{la tripulación} — \textquotedblleft{}the crew\textquotedblright{}.

\begin{quote}
La tripulación es muy amable.
\end{quote}

\begin{quote}
\emph{The crew are very kind.}
\end{quote}

One member is \emph{un tripulante}.

\begin{cousinweb}
\emph{tripulación} is built on \emph{tripular} \textquotedblleft{}to crew a ship\textquotedblright{}, which is usually traced to Latin \emph{interpolāre} \textquotedblleft{}to refurbish, to patch up\textquotedblright{} (EtymDB).
\end{cousinweb}

\begin{grammarlens}[title={What you've built}]
One more word for telling, explaining and describing at B1.
\end{grammarlens}

\subsection*{Your turn}
\begin{itemize}
\raggedright
  \item \emph{Say it:} \emph{la tripulación}
  \item \emph{Say it:} \emph{la tripulación}, once more
  \item \emph{Say it:} a sentence of your own with \emph{la tripulación}
  \item \emph{Recall:} say the Spanish for to land, then the Spanish for the pilot, then say \emph{la tripulación} again
\end{itemize}

\begin{tcolorbox}[breakable,colback=teal!4,colframe=teal!35!black,title={Before you move on}]
What does la tripulación mean? (\textquotedblleft{}The crew\textquotedblright{}.) Say it once more.
\end{tcolorbox}
\section[la azafata]{la azafata — the flight attendant}

\label{lesson:ES-C654-azafata}



\begin{quote}
Before the new one: say the Spanish for the pilot, then the Spanish for the crew.
\end{quote}

\subsection*{You'll want to know: la azafata}
\textbf{la azafata} — \textquotedblleft{}the flight attendant\textquotedblright{}.

\begin{quote}
La azafata nos trae agua.
\end{quote}

\begin{quote}
\emph{The flight attendant brings us water.}
\end{quote}

Latin America says \emph{la aeromoza} or \emph{la sobrecargo}; the neutral term is \emph{auxiliar de vuelo}.

\begin{cousinweb}
\emph{azafata} is built on \emph{azafate} \textquotedblleft{}tray\textquotedblright{}, which comes from Arabic \emph{safaṭ} \textquotedblleft{}basket\textquotedblright{} (EtymDB). The first \emph{azafata} was a lady who held the queen’s tray.
\end{cousinweb}

\begin{grammarlens}[title={What you've built}]
One more word for telling, explaining and describing at B1.
\end{grammarlens}

\subsection*{Your turn}
\begin{itemize}
\raggedright
  \item \emph{Say it:} \emph{la azafata}
  \item \emph{Say it:} \emph{la azafata}, once more
  \item \emph{Say it:} a sentence of your own with \emph{la azafata}
  \item \emph{Recall:} say the Spanish for the pilot, then the Spanish for the crew, then say \emph{la azafata} again
\end{itemize}

\begin{tcolorbox}[breakable,colback=teal!4,colframe=teal!35!black,title={Before you move on}]
What does la azafata mean? (\textquotedblleft{}The flight attendant\textquotedblright{}.) Say it once more.
\end{tcolorbox}
